\documentclass[a4paper, 12pt]{report}

\title{Editeur Web}
%% Pour la police de caractères.
\usepackage{fontspec}
%% Pour la langue des titres et sous-titres.
\usepackage[francais]{babel}
%% Pour de belles images.
\usepackage{graphicx}
%% Pour des marges plus équitables.
\usepackage[margin=2cm]{geometry}
%% Pour faire le glossaire.
\usepackage[xindy]{glossaries}
\makeglossaries
%% La police de caractères.
\setmainfont{Delicious-Roman}
%% Pour pouvoir écrire du code.
%\usepackage{minted}
%\usemintedstyle{tango}
\begin{document}
	\begin{titlepage}
		\center{\includegraphics[width=5cm]{images/logoUM2.png}}	\\ 
		~\\
		~\\
		~\\
		~\\
		~\\		
		\begin{center}
			{\large Rapport préliminaire de projet} \\
			{\large Licence 3}\\
			\vspace{1,5cm}
			{\Huge Editeur de sites web}\\
			~\\
			~\\
			~\\
			\includegraphics[width=12.5cm]{images/logoTest1.png}
			~\\
			~\\
			{\large Réalisé par :} \\
			~\\
			{\LARGE Pierre Burc, Olivier Duplouy, \\
				      Hamza Erraji, Issam Amal,\\
				      Mickaël Berger, Joachim Divet,\\
				      Zidane Sadiki et Abdelhamid Belarbi}\\
			\vspace{1,5cm}
			{\large Sous la direction de :} \\
			~\\
			{\LARGE Michel Meynard} \\
			\vspace{2.5cm}
			{\large Année universitaire 2011-2012}			
		\end{center}
	\end{titlepage}
%%%%%%%%%%%%%%%%%%%%%%%%%%%%%%%%%%%%%%%%%%%%%%%%%%%%%%%%%%%%%%%%%%%%%%%%%%%%%%%%%%%%%%%%%%%%%%%%%%%%%%%%%%%%%%%%%%%%%%%%%%%%%%%%%%
%%%%%%%%%%%%%%%%%%%%%%%%%%%%%%%%%%%%%%%%%%%%%%%%%%%%%%%%%%%%%%%%%%%%%%%%%%%%%%%%%%%%%%%%%%%%%%%%%%%%%%%%%%%%%%%%%%%%%%%%%%%%%%%%%%
\begin{chapter}*{Avant-propos}
	Il serait profondément malsain de commencer la description de ce projet sans, avant cela, une touche de présentation.\\
	Nous sommes donc huit étudiants de la Faculté des sciences de Montpellier, dans le domaine de l'informatique.
	Nous fûmes réunis pour la première fois en janvier 2012 pour débattre sur le sujet de notre projet, à savoir:\\ 
	\quotation{\emph{"Dans le cadre du développement de sites web, on souhaiterait utiliser un éditeur multi-fichiers permettant de réaliser différentes actions sur des fichiers relatifs à un site."}}\\


    Après cette réunion nous avons à force de discussions, d'argumentations opposées, d'idées nouvelles et de feuilles de papier, décidés d'une 
    marche à suivre pour répondre à cette problématique.\\
    Et c'est ainsi que dans les pages qui vont suivre sont exposées nos conclusion quant à la mise en place d'une solution.	
\end{chapter}

	\begin{part}{Conception}
		\begin{chapter}{Analyse de l'existant}
			\begin{section}{Espresso}
				\begin{figure}[h]
					\begin{center}
						\includegraphics[width=3cm]{images/logoEspresso.png}
						\caption{Espresso}
					\end{center}
				\end{figure}~\\
				Le premier programme d'édition de site web que nous avons étudié, il propose les mêmes fonctionnalités (chap. 2) que le
				logiciel que nous nous proposons de développer.\\
				Il sera a priori une bonne source d'inspiration pour nous et lorsque l'on constate a fortiori que l'interface est très soignée, on en déduit que le travail est fait très sérieusement.
			\end{section}
			~\\
			\begin{section}{Aptana}
				\begin{figure}[h]
					\begin{center}
						\includegraphics[width=3cm]{images/logoAptana.png}
						\caption{Aptana}
					\end{center}
				\end{figure}~\\
			\end{section}
			~\\
			\begin{section}{Dreamweaver}
				\begin{figure}[h]
					\begin{center}
						\includegraphics[width=3cm]{images/logoDreamweaver.png}
						\caption{Dreamweaver}
					\end{center}
				\end{figure}~\\
				
			\end{section}
			~\\
		\end{chapter}
		\begin{chapter}{Spécifications}
				
				Implémentation d'un IDE comprenant les fonctionnalités suivantes :

				- édition des fichiers Javascript, PHP, CSS, HTML
				- Un système multi-onglets permettant de naviguer rapidement entre plusieurs fichiers d'un même projet.
				- Possibilité de visualiser le site sous forme d'arborescence.
				- Auto-complétion et coloration syntaxique. 
				- Auto-indentation ( Ou à la demande de l'utilisateur ?).
				- Accès en un clic au manuel PHP/HTML/CSS/Javascript
				- Permettre la validation HTML (Via W3c ?).
				- Bouton ouvrant un naviguateur internet pour visualiser le site.
				- Squelette de site préexistant ???
			~\\
		\end{chapter}
		\begin{chapter}{Modélisation}
		\end{chapter}
	\end{part}
	\begin{part}{Planification}
		\begin{chapter}{Contraintes}
		\end{chapter}
		\begin{chapter}{Choix des outils}

			Pour le développement de l'IDE :

			QTCreator (C++)

			Contrôle de versions du projet :

			GIT

			Pour l'analyse UML :

			yUML.me

			Discussions de groupe :

			Facebook et Micromobs (http://micromobs.com/)

			Diagramme de Gantt :

			www.teamgantt.com




		\end{chapter}
	\end{part}

\end{document}
